\documentclass[uplatex,dvipdfmx]{jsarticle}
\usepackage{amsmath,amssymb}
\usepackage{bm}
\usepackage{float}
\usepackage{mathtools}
\usepackage{ascmac}
\usepackage{mathrsfs}
\usepackage{physics}
\usepackage{here}
\usepackage{url}
\usepackage{slashed}
\usepackage{tikz-cd}
\usepackage[top=25truemm,bottom=30truemm,left=20truemm,right=20truemm]{geometry}
\usepackage{xcolor}
\usepackage{titlesec}
\begin{document}

Fermion場 $\psi$を用いた最もシンプルなbosonicな物理量として
\begin{equation*}
     \mathcal{O}(x) = \bar{\psi}(x)\psi(x) ~~~\mathcal{O}(y) = \bar{\psi}(y)\psi(y)
\end{equation*}
を考える.~Fermionの半交換関係は
\begin{equation*}
    \{\psi(x), \bar{\psi}(y)\} = i\Delta(x-y), \quad \{\psi(x), \psi(y)\} = 0, \quad \{\bar{\psi}(x), \bar{\psi}(y)\} = 0
\end{equation*}
となる.~この半交換関係を用いて,~$\mathcal{O}(x)$ と $\mathcal{O}(y)$ の交換関係を計算すると,
\begin{equation*}
    [AB, CD] = A\{B, C\}D - \{A, C\}BD + C A\{B, D\} - C\{A, D\}B
\end{equation*}
より,~$[\mathcal{O}(x), \mathcal{O}(y)]$ は $\Delta(x-y)$ のような項のみを含むことが分かる.~$x,~y$が空間的に離れているとき,~Lorentz不変性から $\Delta(x-y) = 0$ となることが分かる.~したがって,~
\begin{equation*}
    [\mathcal{O}(x), \mathcal{O}(y)] = 0 \quad \text{for} \quad (x-y)^2 < 0
\end{equation*}
となり,~空間的に離れた場所で独立に観測可能になるので,~因果律が保たれることが分かる.
\end{document}